% !TeX root = ../../Book4.tex
\chapter*{序言}
\addcontentsline{toc}{chapter}{序言}
\markboth{序言}{序言}

取一个整数 \(n\ge2\)。下面这列映射没有什么神秘之处：
\[
0\longrightarrow\mathbb Z\xrightarrow{\times n}\mathbb Z
\longrightarrow\mathbb Z/n\mathbb Z\longrightarrow0.
\]
前一项恰好是后一映射的核，所以它是短正合列。可是，对其中的模取 \(\operatorname{Hom}_{\mathbb Z}(-,\mathbb Z)\) 后，出现的映射 \(\mathbb Z\xrightarrow{\times n}\mathbb Z\) 并不满射；末端留下了 \(\mathbb Z/n\mathbb Z\)。Hom 函子没有保持原来的正合列。本卷将研究这种失效，并计算由此产生的群。

\ProseParagraphBreak
把留下的群记为 \(\operatorname{Ext}^{1}_{\mathbb Z}(\mathbb Z/n\mathbb Z,\mathbb Z)\)，仍需解释为什么更换计算方法不会改变它，以及它怎样随原来的模同态变化。为此要研究复形、同调与消解。张量积失去正合性时出现 Tor，更一般的情形则引出导出函子。这些记号来自具体计算，记录了原函子未能保持的正合关系。

\ProseParagraphBreak
这套方法有一段跨越分支的历史。Poincaré 等人曾借链和边界研究几何空间中不同维度的“洞”；把向量空间换成环上的模以后，同调也能够记录代数映射的核与像之间的差别\cite[Section 1.1]{Weibel1994HomologicalAlgebra}。Eilenberg 与 Mac Lane 在研究自然对应时提出范畴、函子和自然变换的语言；他们1945年的论文题为《General Theory of Natural Equivalences》\cite[Section 7.10]{EtingofEtAl2011RepresentationTheory}。链复形用核与像进行计算，范畴语言则说明这些计算怎样随映射变化；本卷将同时使用这两种方法。

\ProseParagraphBreak
范畴论并不要求忘掉熟悉的对象。给定群、环或模时，我们早已习惯同时研究它们之间的同态；范畴把“对象和保持结构的映射”一起作为讨论对象。自然变换进一步要求一个构造在所有相关映射下相容。比如 \(M\oplus N\) 的两个投影、张量积的泛性质、核与余核的定义，原先各自写出一套公式，本卷第一编将说明它们的共同形式，也说明共同形式之外各自需要什么条件。

\ProseParagraphBreak
第一编从范畴、函子和自然性出发，接着讨论极限、余极限、Yoneda 引理与伴随，再进入加性范畴、交换范畴和图引理。这里的许多结论都要返回模的例子逐项核验：积与直和何时相同，核和余核怎样计算，正合性在哪一步被保持。读者若能在抽象陈述后重新写出一个具体映射，便不容易把泛性质误认为一句只表示“某物存在”的口号。

\ProseParagraphBreak
第二编建立复形和消解，随后构造导出函子、Ext 与 Tor。短正合列为什么会产生同调长正合列，更换投射或内射消解为什么不改变答案，这些都将在相应位置证明。开头的整数例子也会以不同面貌重现：它既可用来计算一个 Ext 群，也能帮助辨认扩张的等价类。后面再讨论万有系数、Künneth 公式与同调维数时，所用的仍是这些基本构造，只是问题换了方向。

\ProseParagraphBreak
第三编把方法放到群上同调、Hochschild 同调和逆极限之中，并建立谱序列。谱序列能把一次难算的同调分成逐步计算，但每一页所得的信息与最终对象之间还有距离；微分、收敛条件和隐藏的扩张问题都须说明。读者第一次遇到它，不妨先盯住一个有界双复形：沿行算、沿列算，各得到什么，为什么最后应当描述同一对象。熟悉这个例子，再读一般构造，会比先背一大片指标更稳妥。

\ProseParagraphBreak
第四编进入同伦范畴和导出范畴。前面为计算而选择的消解，在这里成为比较复形的工具；把拟同构视为可逆之后，可以在新的范畴中讨论导出张量、导出 Hom 与好三角。接下来的 \(t\) 结构又提出一个颇有意思的问题：从保存了各种次数信息的导出范畴中，怎样重新认出类似模的对象？选读的 DG 范畴和导出等价只展开到本卷能够说明的范围，不要求把后续的整套理论压进初次阅读。

\ProseParagraphBreak
本卷主要面向已经学过第二卷模论基础的读者。进入第一编前，应熟悉模同态、商模、张量积和短正合列；投射、内射、平坦模的基本例子会在第二编频繁出现。读本卷不必先学第三卷前半。若正在读第三卷关于深度、局部极小消解和正则局部环的部分，可以先补本卷第五至八章，尤其是维数平移，再回到交换代数中的应用。

\ProseParagraphBreak
阅读顺序也可以随问题调整。想学群上同调，可在消解和导出函子之后读第十一章；想学谱序列，可先读复形、双复形与第十四章，再看它在群扩张或函子复合中的应用；想了解导出范畴，则先掌握同伦、拟同构和有界消解。星号内容可暂放一旁，章节导读会交代后来何处需要它。无论选择哪一段，都请先确认本卷所用的次数、微分和平移方向；同一个熟悉的公式，换一套记号就可能差一个负号。

\ProseParagraphBreak
我希望读者在每个新定义旁留一小块计算的地方。写出一个两项复形，亲自求核、像和同调；把一个短正合列送进 Hom 或张量积，找出哪一端出了问题；声称某个映射“自然”时，画出它应当交换的方块。若计算结果与抽象陈述相符，抽象就有了可以检验的内容。若不符，错误往往也有价值：也许忘了假设，也许把函子的方向弄反了。

\ProseParagraphBreak
从整数模上的一次 Ext 计算，到一个范畴中的导出对象，中间要经过不少定义。它们并不要求读者立刻把每个符号都记住；更要紧的是，始终知道当前想保存哪种关系、失去了什么、又怎样把它找回来。愿这卷书帮助读者在具体计算与一般方法之间自由往返。

\bigskip
\begin{flushright}
R. EverStarine

2026年9月24日
\end{flushright}
